\section{Exact Preprocessing Propositions}
\label{sec:or-preprocessing-proofs}

Fix an identity-closure tagged covering instance after all declared mandatory
columns have been selected. Let $Q$ be the remaining requirement rows,
$A$ the available nonfixed columns, $R_s^Q=R_s\cap Q$, and let each
optional weight be a strictly positive exact rational. There are no
identity-sensitive side constraints, groups, exact-cardinality constraints,
or mutual exclusions: this is the pure covering model. Every transformation
records original identifiers, forced selections, fixed burden, equal-weight
substitutions, and merged-row memberships.

\begin{proposition}[Mandatory carrier and forced propagation]
\label{prop:or-forced-carrier}
If $u\in Q$ has exactly one carrier $s\in A$, every feasible completion
selects $s$. Fixing $s$, adding $w_s$ to the objective offset, deleting
the rows it satisfies, and recomputing carriers preserves feasibility, the
exact optimum value, every optimum identity, and reconstruction.
\end{proposition}

\begin{proof}
Any feasible completion must cover $u$, and $s$ is its only available
carrier, so $s$ occurs in every completion. Removing already satisfied rows
does not change the constraints on the remaining columns. Splitting every
feasible objective into the fixed term $w_s$ and its residual term is a
bijection between feasible residual selections and feasible original
selections containing the already forced columns. Repeating the argument
after every selection proves propagation.
\end{proof}

\begin{proposition}[Coverage dominance]
\label{prop:or-coverage-dominance}
Suppose $s,t\in A$, $s\neq t$, $R_s^Q\subseteq R_t^Q$, and
$w_t\leq w_s$. If $s$ is not mandatory, deleting $s$ preserves
feasibility, the optimum value, and at least one optimum. If $w_t<w_s$, no
optimum contains $s$, so all optimizer identities are preserved. If
$w_t=w_s$, all optimizer identities need not be preserved; a reconstruction
record must retain the admissible substitution $s\leftrightarrow t$ where
its surrounding selection remains feasible.
\end{proposition}

\begin{proof}
Take a feasible selection containing $s$. If it also contains $t$, then
deleting $s$ preserves coverage and strictly lowers burden because
$w_s>0$. If it does not contain $t$, replace $s$ by $t$. Superset
coverage preserves feasibility and the burden weakly falls. Thus every
feasible solution has a no-worse representative without $s$, which proves
feasibility and optimum-value preservation and supplies at least one reduced
optimum. Under strict inequality, either case strictly improves any selection
containing $s$, so no original optimum uses it. Under equality, replacement
can exchange two distinct optimizer identities without changing burden.
\end{proof}

\begin{proposition}[Duplicate columns and empty contribution]
\label{prop:or-duplicate-empty-columns}
For columns with identical residual coverage, retaining a stable
minimum-weight representative and deleting the others preserves feasibility,
the exact optimum value, and at least one optimum. Strictly heavier duplicates
cannot occur in an optimum. Equal-weight duplicates require an explicit
substitution class to reconstruct all optimizer identities. A nonmandatory
positive-weight column with empty residual coverage may be deleted; it occurs
in no optimum.
\end{proposition}

\begin{proof}
Identical-coverage replacement leaves every row status unchanged and weakly
reduces burden, with strict reduction for a heavier duplicate. Stable choice
only fixes which representative remains in the reduced model. Recording the
equal-weight class recovers the alternative identities consistent with the
rest of a lifted selection. An empty-contribution column changes no residual
constraint, and deleting it strictly lowers any selection that contains it.
\end{proof}

\begin{proposition}[Duplicate requirement rows]
\label{prop:or-duplicate-requirements}
If two residual requirement rows have identical carrier sets, replacing them
by one representative preserves feasibility, optimum value, at least one
optimum, and all optimizer identities. Recording both original row identifiers
preserves the feasibility certificate.
\end{proposition}

\begin{proof}
A selection intersects one carrier set if and only if it intersects the
identical other carrier set. The two inequalities therefore define the same
feasible selections and do not enter the objective. Merging them changes no
optimizer; the row-membership record merely expands the final certificate.
\end{proof}

The implementation applies the rules to a stable fixed point and restarts its
ordered scan after every transformation. Each nontrivial step removes a row or
column or fixes a previously available column, so finiteness proves
termination. The complete machine audit distinguishes forced columns,
dominated columns, duplicate columns, empty columns, satisfied rows, and
duplicate rows. These are exact combinatorial propositions. Julia randomized
exhaustive tests and Lean finite-mask identities are separate evidence layers;
neither turns a downstream solver status into a formal optimality proof.
